\section{60--180 天三情景预测}

\subsection{情景设定}

长期情景预测在综合最优参数基础上构建乐观、中性和悲观三种情景。附件真实价格只覆盖到 2026 年 5 月 5 日，之后全部为条件情景外推，不作为真实观测数据。

长期部分不能称为“拟合”：赛题附件只提供到 2026 年 5 月初的真实价格，60--180 天区间没有可用于事后对照的真实观测值。因此本文采用“情景预测”和“平衡点分析”的表述。三情景曲线表示条件中心路径，而不是未来每日收盘价的精确点预测；其可信度由机制一致性、官方外生数据校验、敏感性分析、蒙特卡洛概率区间和状态转移扰动共同评价。

长短期模型并非两套互不相干的模型。短期部分使用同一套物理缺口和市场价格形成递推式，回答封锁初期价格平台如何形成；长期部分仍以该递推机制为中心，只是把未来不可观测的政策协调、绕道恢复、供给响应和二次升级风险从固定参数扩展为情景变量和状态变量。换言之，随着外推期限由 0--60 天延长到 60--180 天，不确定性的主要来源从“参数校准误差”逐步转向“事件状态不确定性”，因此需要用蒙特卡洛情景树和状态转移扰动描述条件区间。

长期模型从短期动态递推模型延伸而来，但在四个方面做了时间尺度调整：

\begin{enumerate}
  \item 供应侧：保持封锁造成的基础缺口，同时允许非海湾产油国和 OPEC 剩余产能缓慢响应。
  \item 缓冲侧：SPR、商业库存和绕道运输继续发挥作用，但其边际能力随时间下降或受运输瓶颈约束。
  \item 需求侧：长期需求价格弹性高于短期弹性，反映高油价持续后消费、生产和替代能源调整更充分。
  \item 风险侧：恐慌溢价逐步衰减，不确定性溢价拆分为前期冲击不确定性和持续封锁制度风险，避免把长期价格简单托在一个常数平台上。
\end{enumerate}

记长期外推期第 $t$ 日有效供应为 $S_t^{(L)}$，有效需求为 $D_t^{(L)}$。在短期模型的基础上，长期供需缺口可写为

\begin{equation}
G_t^{(L)}
=
D_t^{(L)}
-
\left[
S_0-I+R_t+Q_t+B_t+A_t
\right],
\end{equation}

其中 $I$ 为供应中断量，$R_t$ 为实际 SPR 释放量，$Q_t$ 为绕道运输恢复量，$B_t$ 为商业库存缓冲量，$A_t$ 为其他产油国增产或替代供应增量。为避免将 SPR 机械地设为计划上限，本文令实际释放量随缺口和价格压力自适应收缩：

\begin{equation}
R_t=\min\left\{\bar R_t,\ \max\left(G_{t,\mathrm{pre}}+\rho D_0,\ \bar R_t \lambda_t \tau_t\right)\right\},
\end{equation}

其中 $\bar R_t$ 为计划释放能力，$G_{t,\mathrm{pre}}$ 为不含 SPR 时的预释放缺口，$\rho D_0$ 为安全余量，$\lambda_t$ 表示由物理缺口和价格压力共同决定的释放强度，$\tau_t$ 为长期持续后的政策收缩系数。因此，当缺口已基本闭合时，SPR 不会继续按上限满额释放。

长期价格递推采用

\begin{equation}
\hat P_{t+1}
=
(1-v_L)\hat P_t
+v_L
\left[
P_0+\Pi(G_t^{(L)},\varepsilon_L)
\right]
+\Phi_t-\Omega_t+J_t ,
\end{equation}

其中 $\Pi(\cdot)$ 表示由剩余供需缺口和长期需求弹性 $\varepsilon_L$ 共同决定的供需压力；当 $G_t^{(L)}<0$ 时，过剩供给通过折价项向下拉动目标价格。$\Phi_t$ 由地缘风险溢价、前期冲击不确定性和持续封锁制度风险三部分构成，其中制度风险随未解决供应压力和时间信心恢复共同衰减，避免形成长期常数托底。$\Omega_t$ 为政策缓冲确认和预期修复带来的折价，$J_t$ 为库存接近耗尽时的跳变风险项。该形式对应赛题要求中的增产、长期弹性变化、运输瓶颈和库存耗尽风险。

短期和长期风险项的方向不同，并不是前后矛盾。短期窗口内，市场需要完成突发封锁的信息吸收，风险溢价表现为建立和再定价；长期窗口内，市场已经观察到 SPR、库存、绕道和供需调整是否有效，因此风险项应随未解决缺口和信心恢复而逐步衰减。换言之，短期模型回答“冲击如何形成价格平台”，长期模型回答“平台在政策缓冲和需求调整后如何回归，尾部风险在哪里重新抬升”。

\begin{table}[htbp]
\centering
\caption{长期机制强度参数及检验边界}
\begin{tabularx}{\linewidth}{lrrX}
\toprule
机制参数 & 基准值 & 扰动范围 & 含义与边界 \\
\midrule
SPR 收缩起点 & 75 & 60--105 & SPR 从应急计划释放转向缺口驱动收缩的开始时间 \\
冲击不确定性占比 & 0.45 & 0.25--0.65 & 前期冲击再定价在长期不确定性中的贡献，不作为外部官方统计量 \\
制度风险占比 & 0.90 & 0.50--1.30 & 持续封锁状态对长期风险溢价的贡献强度 \\
信心恢复后风险底座 & 0.40 & 0.20--0.60 & 缓冲机制被市场确认后仍保留的制度风险比例 \\
过剩供给回归强度 & 1.35 & 0.80--2.20 & 供给超过需求时向下修正目标价格的机制强度 \\
封锁风险衰减速度 & 0.004 & 0.002--0.008 & 封锁风险溢价随时间被市场吸收的日度速度 \\
\bottomrule
\end{tabularx}
\end{table}

\subsection{长期参数来源与校验边界}

长期模型中的部分参数不是可以从附件直接读取的真实观测量，而是用于描述恢复速度、政策收缩和制度风险衰减的机制参数。为避免参数自由度过高，本文对长期参数采用三条约束：一是方向必须来自赛题机制或能源市场逻辑；二是数值只能作为基准情景假设进入，不能冒充官方统计数据；三是所有关键参数必须进入扰动区间、敏感性排序或蒙特卡洛检验。长期参数来源与校验方式见表 \ref{tab:long-parameter-audit}。

\begin{table}[htbp]
\centering
\caption{长期关键参数来源与校验方式}
\label{tab:long-parameter-audit}
\begin{tabularx}{\linewidth}{lX X}
\toprule
参数组 & 来源逻辑 & 校验方式 \\
\midrule
SPR 收缩起点与释放强度 & 来自赛题给出的 SPR 缓冲机制，并结合“长期封锁下不能无限满额释放”的政策约束。 & 在 60--105 天范围内扰动，检验其对第 180 天价格和尾部峰值的影响。 \\
冲击不确定性与制度风险 & 来自封锁未解除时保险、航运、二次升级和政策协调失败的持续风险。 & 与 GPR 滞后检验、蒙特卡洛情景树和状态转移情景树合读，不作为同步价格拟合因子。 \\
封锁风险衰减速度 & 表示市场对封锁信息逐步吸收和政策缓冲被确认后的风险消化速度。 & 设置 0.002--0.008 日衰减率区间，并在敏感性分析中单独排序，避免隐藏常数。 \\
过剩供给回归强度 & 表示当缓冲、增产和需求调整使供给压力缓解后，目标价格向下修正的速度。 & 通过扰动区间和情景曲线检验，防止中性路径被人为拉成固定直线。 \\
其他产油国增产与绕道恢复 & 来自赛题要求的增产、运输瓶颈和绕道能力设定，并用 JODI 多国月度产量、进出口和炼厂产出作数量级校验。 & 只用于长期条件外推，不反向用于短期拟合；JODI 为多国上报口径，不能写成完整全球总量。 \\
\bottomrule
\end{tabularx}
\end{table}

\keypoint{长期参数的角色是描述机制强弱和恢复速度，而不是替代真实数据。本文把这些参数显式列出、放入扰动区间，并用敏感性和蒙特卡洛结果报告其不确定性。}

\begin{table}[htbp]
\centering
\caption{三情景关键参数}
\begin{tabular}{lrrrrrr}
\toprule
情景 & 供应中断 & SPR 上限 & 绕道能力 & 长期弹性 & 恐慌衰减 & 风险权重 \\
\midrule
乐观 & 1400 & 700 & 300 & -0.25 & 0.12 & 2.02 \\
中性 & 1493 & 670 & 237 & -0.25 & 0.11 & 2.46 \\
悲观 & 1800 & 200 & 150 & -0.10 & 0.04 & 3.32 \\
\bottomrule
\end{tabular}
\end{table}

三情景参数并非三组孤立猜测。乐观情景代表高 SPR 可用量、绕道恢复较充分、长期需求调整较快；悲观情景代表供应中断接近上界、SPR 释放受限、绕道能力偏低且制度风险溢价持续；中性情景则采用综合最优参数附近的中心设定。后续敏感性分析和蒙特卡洛抽样用于检验：当这些假设联合变化时，价格中心路径和尾部风险是否仍保持一致。

在补充官方外生数据后，本文将 EIA、JODI 与 OPEC 从“论文证据层”推进到“长期参数约束层”。具体做法不是用外部数据替代赛题参数，而是由 JODI 多国同口径产量、库存、进出口、需求、炼厂产出、EIA 美国 SPR 和商业库存变化，以及 OPEC 全球供需平衡表中的需求增量、非 DoC 供给增量和 DoC 原油需求差额生成温和乘数，对绕道能力、库存缓冲、需求收缩、长期弹性、SPR 上限可信度和风险权重进行约束。得到的乘数见表 \ref{tab:official-factor-multipliers}。

\begin{table}[htbp]
\centering
\caption{官方外生数据生成的长期参数约束因子}
\label{tab:official-factor-multipliers}
\begin{tabularx}{\linewidth}{l r X}
\toprule
约束因子 & 乘数 & 解释 \\
\midrule
绕道和替代贸易恢复能力 & 0.9725 & JODI 进出口同口径收缩，说明替代贸易恢复不宜过度乐观 \\
商业库存日缓冲上限 & 1.0739 & JODI 原油/成品油库存和 EIA 商业库存均提供缓冲数量级支撑 \\
长期需求收缩幅度 & 0.9560 & JODI 与 OPEC 均显示需求韧性，需求破坏不宜设得过强 \\
长期需求弹性绝对值 & 0.9750 & 高价下需求会调整，但 OPEC 正常需求增长使长期弹性略收敛 \\
SPR 释放上限可信度 & 0.9672 & EIA 美国 SPR 实际释放较温和，因此全球协调释放上限作小幅折扣 \\
贸易收缩风险权重 & 1.0543 & JODI 进出口收缩叠加 OPEC DoC 需求差额上升，抬升尾部风险 \\
长期不确定性强度 & 1.0303 & 贸易收缩和下半年需求差额峰值使制度风险略上调 \\
\bottomrule
\end{tabularx}
\end{table}

\subsection{预测结果}

\begin{table}[htbp]
\centering
\caption{60--180 天三情景预测结果}
\begin{tabular}{lrrrrrrr}
\toprule
情景 & 第60天 & 第90天 & 第120天 & 第180天 & 外推期最高价 & 外推期均价 & 风险 \\
\midrule
乐观 & 109.77 & 94.36 & 92.79 & 87.78 & 103.79 & 92.16 & 低 \\
中性 & 109.77 & 107.12 & 102.43 & 94.15 & 110.68 & 102.17 & 低 \\
悲观 & 109.77 & 127.81 & 118.64 & 110.16 & 133.91 & 120.01 & 高 \\
\bottomrule
\end{tabular}
\end{table}

\begin{figure}[htbp]
\centering
\includegraphics[width=0.98\linewidth]{蒙特卡洛情景树高级组合图.png}
\caption{长期情景预测主图：蒙特卡洛情景树高级组合图。左侧展示中心路径、抽样路径云和 5\%--95\% / 25\%--75\% 条件区间，右上展示第 180 天终点分布，右下展示外推期突破 120、130 美元/桶的累计概率曲线；该图证明长期结论应读作条件区间和尾部概率，而不是单点预测。}
\label{fig:long-term-fan}
\end{figure}

\begin{figure}[htbp]
\centering
\includegraphics[width=0.94\linewidth]{短期拟合与长期预测总览.png}
\caption{短期机制拟合与 60--180 天三情景预测总览。竖线左侧为附件真实数据与短期模型拟合，竖线右侧为情景外推，不作为真实观测。}
\end{figure}

图 \ref{fig:long-term-fan} 给出长期情景预测的主图。为刻画中心路径之外的事件状态不确定性，本文额外建立状态转移扰动层。该层不改变三情景物理中心路径，而是让未来状态在缓和、维持和升级之间随机切换：缓和状态向乐观中心靠近，维持状态围绕中性中心，升级状态向悲观中心和跳涨方向移动；扰动强度由附件历史价格波动估计后进行保守折扣。进一步地，本文用 GPR 事件月历史分位数、OVX 冲突窗口历史分位数以及滞后 OVX 与实现波动率相关性生成综合状态转移压力，用于温和提高早中期升级概率。这样得到的长期图不再是一条孤立直线，而是中心路径、状态概率和尾部区间的组合。

\begin{figure}[htbp]
\centering
\includegraphics[width=0.94\linewidth]{长期状态转移情景树.png}
\caption{长期状态转移情景树。上图为状态转移后的价格扇形区间，下图为缓和、维持、升级三类状态的样本占比；该图用于回应长期中性线过平的问题，并展示二次升级风险如何进入尾部区间。}
\label{fig:state-transition-fan}
\end{figure}

结果显示，中性情景下价格在 180 天附近回落至约 94.15 美元/桶；悲观情景下，若供应中断偏大、SPR 释放弱、绕道恢复慢且风险溢价维持高位，价格仍可能再次升至 120 美元/桶以上，并在外推期达到约 133.91 美元/桶的高点。图 \ref{fig:state-transition-fan} 的状态转移增强模型进一步给出第 180 天 5\%--95\% 条件区间 86.36--96.06 美元/桶，外推期最高价 P95 为 130.51 美元/桶，突破 120 美元/桶的条件概率为 38.3\%，说明长期中心路径的日度精度不能单独夸大，真正应关注的是尾部区间和状态切换风险。

\begin{table}[htbp]
\centering
\caption{长期情景预测的可信度评价框架}
\begin{tabularx}{\linewidth}{lX X}
\toprule
评价维度 & 本文证据 & 解释边界 \\
\midrule
机制一致性 & SPR 自适应收缩、长期需求弹性、供给恢复和风险溢价衰减进入同一递推式 & 检验情景逻辑是否自洽，不代表未来真实价格已经被验证 \\
外部数量级校验 & EIA 周度 SPR、商业库存和美国原油产量，叠加 JODI 多国月度产量、库存、进出口、需求和炼厂产出、OPEC 全球供需平衡表，以及 OVX 隐含波动率 & EIA 为美国口径，JODI 为多国上报口径，OPEC 为正常基线预测，OVX 为期权市场波动预期；这些数据用于数量级和风险约束，不能直接替代赛题冲击参数 \\
单因素敏感性 & 第 180 天价格对封锁风险衰减速度、地缘风险权重和制度风险强度最敏感 & 识别关键变量，不给出联合概率 \\
蒙特卡洛情景树 & 第 180 天价格 P05--P95 为 89.58--106.44 美元/桶，最高价突破 120 概率约 16.3\% & 条件概率，不是无条件市场概率 \\
状态转移情景树 & 第 180 天价格 P05--P95 为 86.36--96.06 美元/桶，外推期最高价 P95 为 130.51 美元/桶，突破 120 概率约 38.3\% & 用于解释平滑中心线之外的事件状态切换和短期扰动，不代表未来真实价格已被观测 \\
\bottomrule
\end{tabularx}
\end{table}

从价格平衡点看，乐观情景最终回落至约 87.78 美元/桶，说明若 SPR 可用上限较高、绕道运输恢复充分且长期需求弹性发挥作用，油价可以逐步脱离 110--120 美元平台。中性情景第 180 天约为 94.15 美元/桶，更接近政策缓冲与市场风险重新平衡后的中心水平；此时实际 SPR 释放已经从计划上限自动收缩，且物理缺口基本闭合，说明价格回落不是依赖政策变量机械满额托底。悲观情景第 180 天仍为 110.16 美元/桶，说明若供应中断规模接近上界、SPR 释放不足且制度风险溢价持续，市场仍可能维持相对高位，并保留二次跳涨尾部风险。

需要主动说明的是，中性第 180 天物理缺口已基本闭合，但价格仍高于战前 75 美元/桶，并非供需项失效，而是模型中仍保留不确定性和制度风险溢价。该设定表示：即便物理缺口基本闭合，只要霍尔木兹封锁的制度状态没有完全解除，市场仍会对运输保险、航运安全、库存再补充和再次升级风险给出残余溢价。因此 94.15 美元/桶不是模型声称的唯一真实价格，而是“物理缺口闭合但制度风险尚未完全消失”条件下的中心路径。

\begin{figure}[htbp]
\centering
\includegraphics[width=0.92\linewidth]{库存与供需缺口风险.png}
\caption{商业库存剩余量与剩余供需缺口变化。悲观情景下缓冲能力下降、剩余缺口更难闭合，因此二次跳涨风险主要来自库存约束与制度风险叠加。}
\end{figure}
